% Options for packages loaded elsewhere
\PassOptionsToPackage{unicode}{hyperref}
\PassOptionsToPackage{hyphens}{url}
\PassOptionsToPackage{dvipsnames,svgnames,x11names}{xcolor}
%
\documentclass[
  letterpaper,
  DIV=11,
  numbers=noendperiod]{scrartcl}

\usepackage{amsmath,amssymb}
\usepackage{iftex}
\ifPDFTeX
  \usepackage[T1]{fontenc}
  \usepackage[utf8]{inputenc}
  \usepackage{textcomp} % provide euro and other symbols
\else % if luatex or xetex
  \usepackage{unicode-math}
  \defaultfontfeatures{Scale=MatchLowercase}
  \defaultfontfeatures[\rmfamily]{Ligatures=TeX,Scale=1}
\fi
\usepackage{lmodern}
\ifPDFTeX\else  
    % xetex/luatex font selection
\fi
% Use upquote if available, for straight quotes in verbatim environments
\IfFileExists{upquote.sty}{\usepackage{upquote}}{}
\IfFileExists{microtype.sty}{% use microtype if available
  \usepackage[]{microtype}
  \UseMicrotypeSet[protrusion]{basicmath} % disable protrusion for tt fonts
}{}
\makeatletter
\@ifundefined{KOMAClassName}{% if non-KOMA class
  \IfFileExists{parskip.sty}{%
    \usepackage{parskip}
  }{% else
    \setlength{\parindent}{0pt}
    \setlength{\parskip}{6pt plus 2pt minus 1pt}}
}{% if KOMA class
  \KOMAoptions{parskip=half}}
\makeatother
\usepackage{xcolor}
\setlength{\emergencystretch}{3em} % prevent overfull lines
\setcounter{secnumdepth}{-\maxdimen} % remove section numbering
% Make \paragraph and \subparagraph free-standing
\makeatletter
\ifx\paragraph\undefined\else
  \let\oldparagraph\paragraph
  \renewcommand{\paragraph}{
    \@ifstar
      \xxxParagraphStar
      \xxxParagraphNoStar
  }
  \newcommand{\xxxParagraphStar}[1]{\oldparagraph*{#1}\mbox{}}
  \newcommand{\xxxParagraphNoStar}[1]{\oldparagraph{#1}\mbox{}}
\fi
\ifx\subparagraph\undefined\else
  \let\oldsubparagraph\subparagraph
  \renewcommand{\subparagraph}{
    \@ifstar
      \xxxSubParagraphStar
      \xxxSubParagraphNoStar
  }
  \newcommand{\xxxSubParagraphStar}[1]{\oldsubparagraph*{#1}\mbox{}}
  \newcommand{\xxxSubParagraphNoStar}[1]{\oldsubparagraph{#1}\mbox{}}
\fi
\makeatother

\usepackage{color}
\usepackage{fancyvrb}
\newcommand{\VerbBar}{|}
\newcommand{\VERB}{\Verb[commandchars=\\\{\}]}
\DefineVerbatimEnvironment{Highlighting}{Verbatim}{commandchars=\\\{\}}
% Add ',fontsize=\small' for more characters per line
\usepackage{framed}
\definecolor{shadecolor}{RGB}{241,243,245}
\newenvironment{Shaded}{\begin{snugshade}}{\end{snugshade}}
\newcommand{\AlertTok}[1]{\textcolor[rgb]{0.68,0.00,0.00}{#1}}
\newcommand{\AnnotationTok}[1]{\textcolor[rgb]{0.37,0.37,0.37}{#1}}
\newcommand{\AttributeTok}[1]{\textcolor[rgb]{0.40,0.45,0.13}{#1}}
\newcommand{\BaseNTok}[1]{\textcolor[rgb]{0.68,0.00,0.00}{#1}}
\newcommand{\BuiltInTok}[1]{\textcolor[rgb]{0.00,0.23,0.31}{#1}}
\newcommand{\CharTok}[1]{\textcolor[rgb]{0.13,0.47,0.30}{#1}}
\newcommand{\CommentTok}[1]{\textcolor[rgb]{0.37,0.37,0.37}{#1}}
\newcommand{\CommentVarTok}[1]{\textcolor[rgb]{0.37,0.37,0.37}{\textit{#1}}}
\newcommand{\ConstantTok}[1]{\textcolor[rgb]{0.56,0.35,0.01}{#1}}
\newcommand{\ControlFlowTok}[1]{\textcolor[rgb]{0.00,0.23,0.31}{\textbf{#1}}}
\newcommand{\DataTypeTok}[1]{\textcolor[rgb]{0.68,0.00,0.00}{#1}}
\newcommand{\DecValTok}[1]{\textcolor[rgb]{0.68,0.00,0.00}{#1}}
\newcommand{\DocumentationTok}[1]{\textcolor[rgb]{0.37,0.37,0.37}{\textit{#1}}}
\newcommand{\ErrorTok}[1]{\textcolor[rgb]{0.68,0.00,0.00}{#1}}
\newcommand{\ExtensionTok}[1]{\textcolor[rgb]{0.00,0.23,0.31}{#1}}
\newcommand{\FloatTok}[1]{\textcolor[rgb]{0.68,0.00,0.00}{#1}}
\newcommand{\FunctionTok}[1]{\textcolor[rgb]{0.28,0.35,0.67}{#1}}
\newcommand{\ImportTok}[1]{\textcolor[rgb]{0.00,0.46,0.62}{#1}}
\newcommand{\InformationTok}[1]{\textcolor[rgb]{0.37,0.37,0.37}{#1}}
\newcommand{\KeywordTok}[1]{\textcolor[rgb]{0.00,0.23,0.31}{\textbf{#1}}}
\newcommand{\NormalTok}[1]{\textcolor[rgb]{0.00,0.23,0.31}{#1}}
\newcommand{\OperatorTok}[1]{\textcolor[rgb]{0.37,0.37,0.37}{#1}}
\newcommand{\OtherTok}[1]{\textcolor[rgb]{0.00,0.23,0.31}{#1}}
\newcommand{\PreprocessorTok}[1]{\textcolor[rgb]{0.68,0.00,0.00}{#1}}
\newcommand{\RegionMarkerTok}[1]{\textcolor[rgb]{0.00,0.23,0.31}{#1}}
\newcommand{\SpecialCharTok}[1]{\textcolor[rgb]{0.37,0.37,0.37}{#1}}
\newcommand{\SpecialStringTok}[1]{\textcolor[rgb]{0.13,0.47,0.30}{#1}}
\newcommand{\StringTok}[1]{\textcolor[rgb]{0.13,0.47,0.30}{#1}}
\newcommand{\VariableTok}[1]{\textcolor[rgb]{0.07,0.07,0.07}{#1}}
\newcommand{\VerbatimStringTok}[1]{\textcolor[rgb]{0.13,0.47,0.30}{#1}}
\newcommand{\WarningTok}[1]{\textcolor[rgb]{0.37,0.37,0.37}{\textit{#1}}}

\providecommand{\tightlist}{%
  \setlength{\itemsep}{0pt}\setlength{\parskip}{0pt}}\usepackage{longtable,booktabs,array}
\usepackage{calc} % for calculating minipage widths
% Correct order of tables after \paragraph or \subparagraph
\usepackage{etoolbox}
\makeatletter
\patchcmd\longtable{\par}{\if@noskipsec\mbox{}\fi\par}{}{}
\makeatother
% Allow footnotes in longtable head/foot
\IfFileExists{footnotehyper.sty}{\usepackage{footnotehyper}}{\usepackage{footnote}}
\makesavenoteenv{longtable}
\usepackage{graphicx}
\makeatletter
\def\maxwidth{\ifdim\Gin@nat@width>\linewidth\linewidth\else\Gin@nat@width\fi}
\def\maxheight{\ifdim\Gin@nat@height>\textheight\textheight\else\Gin@nat@height\fi}
\makeatother
% Scale images if necessary, so that they will not overflow the page
% margins by default, and it is still possible to overwrite the defaults
% using explicit options in \includegraphics[width, height, ...]{}
\setkeys{Gin}{width=\maxwidth,height=\maxheight,keepaspectratio}
% Set default figure placement to htbp
\makeatletter
\def\fps@figure{htbp}
\makeatother
% definitions for citeproc citations
\NewDocumentCommand\citeproctext{}{}
\NewDocumentCommand\citeproc{mm}{%
  \begingroup\def\citeproctext{#2}\cite{#1}\endgroup}
\makeatletter
 % allow citations to break across lines
 \let\@cite@ofmt\@firstofone
 % avoid brackets around text for \cite:
 \def\@biblabel#1{}
 \def\@cite#1#2{{#1\if@tempswa , #2\fi}}
\makeatother
\newlength{\cslhangindent}
\setlength{\cslhangindent}{1.5em}
\newlength{\csllabelwidth}
\setlength{\csllabelwidth}{3em}
\newenvironment{CSLReferences}[2] % #1 hanging-indent, #2 entry-spacing
 {\begin{list}{}{%
  \setlength{\itemindent}{0pt}
  \setlength{\leftmargin}{0pt}
  \setlength{\parsep}{0pt}
  % turn on hanging indent if param 1 is 1
  \ifodd #1
   \setlength{\leftmargin}{\cslhangindent}
   \setlength{\itemindent}{-1\cslhangindent}
  \fi
  % set entry spacing
  \setlength{\itemsep}{#2\baselineskip}}}
 {\end{list}}
\usepackage{calc}
\newcommand{\CSLBlock}[1]{\hfill\break\parbox[t]{\linewidth}{\strut\ignorespaces#1\strut}}
\newcommand{\CSLLeftMargin}[1]{\parbox[t]{\csllabelwidth}{\strut#1\strut}}
\newcommand{\CSLRightInline}[1]{\parbox[t]{\linewidth - \csllabelwidth}{\strut#1\strut}}
\newcommand{\CSLIndent}[1]{\hspace{\cslhangindent}#1}

\KOMAoption{captions}{tableheading}
\usepackage{amsmath}
\usepackage{amssymb}
\makeatletter
\@ifpackageloaded{tcolorbox}{}{\usepackage[skins,breakable]{tcolorbox}}
\@ifpackageloaded{fontawesome5}{}{\usepackage{fontawesome5}}
\definecolor{quarto-callout-color}{HTML}{909090}
\definecolor{quarto-callout-note-color}{HTML}{0758E5}
\definecolor{quarto-callout-important-color}{HTML}{CC1914}
\definecolor{quarto-callout-warning-color}{HTML}{EB9113}
\definecolor{quarto-callout-tip-color}{HTML}{00A047}
\definecolor{quarto-callout-caution-color}{HTML}{FC5300}
\definecolor{quarto-callout-color-frame}{HTML}{acacac}
\definecolor{quarto-callout-note-color-frame}{HTML}{4582ec}
\definecolor{quarto-callout-important-color-frame}{HTML}{d9534f}
\definecolor{quarto-callout-warning-color-frame}{HTML}{f0ad4e}
\definecolor{quarto-callout-tip-color-frame}{HTML}{02b875}
\definecolor{quarto-callout-caution-color-frame}{HTML}{fd7e14}
\makeatother
\makeatletter
\@ifpackageloaded{caption}{}{\usepackage{caption}}
\AtBeginDocument{%
\ifdefined\contentsname
  \renewcommand*\contentsname{Table of contents}
\else
  \newcommand\contentsname{Table of contents}
\fi
\ifdefined\listfigurename
  \renewcommand*\listfigurename{List of Figures}
\else
  \newcommand\listfigurename{List of Figures}
\fi
\ifdefined\listtablename
  \renewcommand*\listtablename{List of Tables}
\else
  \newcommand\listtablename{List of Tables}
\fi
\ifdefined\figurename
  \renewcommand*\figurename{Figure}
\else
  \newcommand\figurename{Figure}
\fi
\ifdefined\tablename
  \renewcommand*\tablename{Table}
\else
  \newcommand\tablename{Table}
\fi
}
\@ifpackageloaded{float}{}{\usepackage{float}}
\floatstyle{ruled}
\@ifundefined{c@chapter}{\newfloat{codelisting}{h}{lop}}{\newfloat{codelisting}{h}{lop}[chapter]}
\floatname{codelisting}{Listing}
\newcommand*\listoflistings{\listof{codelisting}{List of Listings}}
\makeatother
\makeatletter
\makeatother
\makeatletter
\@ifpackageloaded{caption}{}{\usepackage{caption}}
\@ifpackageloaded{subcaption}{}{\usepackage{subcaption}}
\makeatother

\ifLuaTeX
  \usepackage{selnolig}  % disable illegal ligatures
\fi
\usepackage{bookmark}

\IfFileExists{xurl.sty}{\usepackage{xurl}}{} % add URL line breaks if available
\urlstyle{same} % disable monospaced font for URLs
\hypersetup{
  pdftitle={🚧Experience-Weighted Attraction},
  colorlinks=true,
  linkcolor={blue},
  filecolor={Maroon},
  citecolor={Blue},
  urlcolor={Blue},
  pdfcreator={LaTeX via pandoc}}


\title{🚧Experience-Weighted Attraction}
\author{}
\date{}

\begin{document}
\maketitle


\subsection{General Principles}\label{general-principles}

Social learning models often aim to disentangle asocial learning
(learning from personal experience) from social learning (learning from
others). The \textbf{Experience-Weighted Attraction (EWA)} model
Barrett, McElreath, and Perry (2017) provides a robust framework for
this. In the context of the Panama capuchin monkey study Barrett,
McElreath, and Perry (2017), the model tracks how individuals update
their ``attractions'' to different foraging techniques based on their
own success and the social cues they observe from others.

The core idea is that individuals have a set of attractions
\(A_{i,j,t}\) for each technique \(j\). These attractions are updated
over time, and the probability of choosing a technique is a mixture of
asocial preferences (derived from attractions) and social influence.

In this documentation, we present two primary variations of the EWA
model adapted from Barrett, McElreath, and Perry (2017):

\begin{enumerate}
\def\labelenumi{\arabic{enumi}.}
\item
  \textbf{Social Learning (EWA)}: A streamlined version (Time Frequency
  Pay-off model) focusing on the interplay between individual experience
  (asocial learning) and social frequency biased by payoffs. It uses a
  reduced parameter set (4 varying effects) to identify the core drivers
  of behavioral diffusion.
\item
  \textbf{Social Learning (EWA-ILV)}: A comprehensive hierarchical model
  (Experience-Weighted Attraction with Individual Level Covariates) that
  evaluates multiple social biases (e.g., kinship, prestige, age
  similarity) and explicitly models how an individual's age influences
  their learning rate (\(\phi\)) and reliance on social cues
  (\(\gamma\)).
\end{enumerate}

\begin{center}\rule{0.5\linewidth}{0.5pt}\end{center}

\subsection{1. Social Learning (EWA)}\label{social-learning-ewa}

The ``EWA'' model is a basic version that includes asocial learning and
pay-off biased social learning. It evaluates four varying effects:
learning rate, social weight, conformity, and pay-off bias.

\subsubsection{Example}\label{example}

\begin{tcolorbox}[enhanced jigsaw, colback=white, toprule=.15mm, coltitle=black, colframe=quarto-callout-note-color-frame, breakable, opacityback=0, bottomrule=.15mm, leftrule=.75mm, opacitybacktitle=0.6, colbacktitle=quarto-callout-note-color!10!white, title=\textcolor{quarto-callout-note-color}{\faInfo}\hspace{0.5em}{Note}, left=2mm, bottomtitle=1mm, toptitle=1mm, titlerule=0mm, arc=.35mm, rightrule=.15mm]

\begin{itemize}
\tightlist
\item
  \textbf{Hierarchical Structure}: Intercepts for learning rates and
  social influence vary by individual (\texttt{id}). Barrett, McElreath,
  and Perry (2017) use a Cholesky-factored multivariate normal prior to
  account for correlations between these effects.
\item
  \textbf{Dtype Sensitivity}: Ensure all data arrays are
  \texttt{float64} to match JAX's 64-bit backend requirements for
  numerical stability.
\end{itemize}

\end{tcolorbox}

\subsection{Python}

\begin{Shaded}
\begin{Highlighting}[]
\ImportTok{from}\NormalTok{ BI }\ImportTok{import}\NormalTok{ bi}
\ImportTok{import}\NormalTok{ jax.numpy }\ImportTok{as}\NormalTok{ jnp}
\ImportTok{import}\NormalTok{ jax}

\NormalTok{m }\OperatorTok{=}\NormalTok{ bi(platform}\OperatorTok{=}\StringTok{\textquotesingle{}cpu\textquotesingle{}}\NormalTok{)}

\KeywordTok{def}\NormalTok{ model(K, J, tech, }\BuiltInTok{id}\NormalTok{, bout, y\_prev, s, ps):}
    \CommentTok{\# Priors}
\NormalTok{    lam   }\OperatorTok{=}\NormalTok{ m.dist.exponential(name}\OperatorTok{=}\StringTok{"lambda"}\NormalTok{, rate}\OperatorTok{=}\FloatTok{1.0}\NormalTok{)}
\NormalTok{    mu    }\OperatorTok{=}\NormalTok{ m.dist.normal(name}\OperatorTok{=}\StringTok{"mu"}\NormalTok{, loc}\OperatorTok{=}\NormalTok{jnp.zeros(}\DecValTok{4}\NormalTok{), scale}\OperatorTok{=}\FloatTok{1.0}\NormalTok{)}
\NormalTok{    sigma }\OperatorTok{=}\NormalTok{ m.dist.exponential(name}\OperatorTok{=}\StringTok{"sigma"}\NormalTok{, rate}\OperatorTok{=}\NormalTok{jnp.full(}\DecValTok{4}\NormalTok{, }\FloatTok{3.0}\NormalTok{))}
\NormalTok{    L\_Rho }\OperatorTok{=}\NormalTok{ m.dist.lkj\_cholesky(name}\OperatorTok{=}\StringTok{"L\_Rho"}\NormalTok{, dimension}\OperatorTok{=}\DecValTok{4}\NormalTok{, concentration}\OperatorTok{=}\FloatTok{3.0}\NormalTok{)}
\NormalTok{    z     }\OperatorTok{=}\NormalTok{ m.dist.normal(name}\OperatorTok{=}\StringTok{"z"}\NormalTok{, loc}\OperatorTok{=}\NormalTok{jnp.zeros((}\DecValTok{4}\NormalTok{, J)), scale}\OperatorTok{=}\FloatTok{1.0}\NormalTok{)}

    \CommentTok{\# Varying effects}
\NormalTok{    L\_scaled }\OperatorTok{=}\NormalTok{ L\_Rho }\OperatorTok{*}\NormalTok{ sigma[:, }\VariableTok{None}\NormalTok{]}
\NormalTok{    a\_id     }\OperatorTok{=}\NormalTok{ (L\_scaled }\OperatorTok{@}\NormalTok{ z).T}

    \CommentTok{\# Parameters per individual}
\NormalTok{    phi\_arr   }\OperatorTok{=}\NormalTok{ jax.nn.sigmoid(mu[}\DecValTok{0}\NormalTok{] }\OperatorTok{+}\NormalTok{ a\_id[}\BuiltInTok{id}\NormalTok{, }\DecValTok{0}\NormalTok{])      }\CommentTok{\# Learning rate}
\NormalTok{    gamma\_arr }\OperatorTok{=}\NormalTok{ jax.nn.sigmoid(mu[}\DecValTok{1}\NormalTok{] }\OperatorTok{+}\NormalTok{ a\_id[}\BuiltInTok{id}\NormalTok{, }\DecValTok{1}\NormalTok{])      }\CommentTok{\# Social weight}
\NormalTok{    fconf\_arr }\OperatorTok{=}\NormalTok{ jnp.exp(mu[}\DecValTok{2}\NormalTok{] }\OperatorTok{+}\NormalTok{ a\_id[}\BuiltInTok{id}\NormalTok{, }\DecValTok{2}\NormalTok{])             }\CommentTok{\# Conformity}
\NormalTok{    Bpay\_arr  }\OperatorTok{=}\NormalTok{ mu[}\DecValTok{3}\NormalTok{] }\OperatorTok{+}\NormalTok{ a\_id[}\BuiltInTok{id}\NormalTok{, }\DecValTok{3}\NormalTok{]                      }\CommentTok{\# Pay{-}off bias}

    \CommentTok{\# Scan carry: Attractions}
    \KeywordTok{def}\NormalTok{ step\_fn(carry, x):}
\NormalTok{        AC }\OperatorTok{=}\NormalTok{ carry}
\NormalTok{        (phi\_i, gamma\_i, fconf\_i, Bpay\_i, tech\_i, id\_i, bout\_i, prev\_y\_i, s\_i, ps\_i) }\OperatorTok{=}\NormalTok{ x}

        \CommentTok{\# Update attractions (EWA)}
\NormalTok{        ac\_new }\OperatorTok{=}\NormalTok{ jnp.where(bout\_i }\OperatorTok{==} \DecValTok{1}\NormalTok{, jnp.zeros(K), (}\FloatTok{1.0} \OperatorTok{{-}}\NormalTok{ phi\_i) }\OperatorTok{*}\NormalTok{ AC[id\_i] }\OperatorTok{+}\NormalTok{ phi\_i }\OperatorTok{*}\NormalTok{ prev\_y\_i)}
        
        \CommentTok{\# Asocial choice probability}
\NormalTok{        prob\_individual }\OperatorTok{=}\NormalTok{ jax.nn.softmax(lam }\OperatorTok{*}\NormalTok{ ac\_new)}

        \CommentTok{\# Social choice probability (Pay{-}off bias)}
\NormalTok{        lin\_mod }\OperatorTok{=}\NormalTok{ jnp.concatenate([jnp.ones(}\DecValTok{1}\NormalTok{), jnp.exp(Bpay\_i }\OperatorTok{*}\NormalTok{ ps\_i[}\DecValTok{1}\NormalTok{:])]) }\OperatorTok{*}\NormalTok{ jnp.power(s\_i, fconf\_i)}
\NormalTok{        safe\_denom }\OperatorTok{=}\NormalTok{ jnp.where(jnp.}\BuiltInTok{sum}\NormalTok{(lin\_mod) }\OperatorTok{\textgreater{}} \DecValTok{0}\NormalTok{, jnp.}\BuiltInTok{sum}\NormalTok{(lin\_mod), }\FloatTok{1.0}\NormalTok{)}
\NormalTok{        prob\_social }\OperatorTok{=}\NormalTok{ lin\_mod }\OperatorTok{/}\NormalTok{ safe\_denom}

        \CommentTok{\# Mixture probability}
\NormalTok{        prob\_mix }\OperatorTok{=}\NormalTok{ (}\FloatTok{1.0} \OperatorTok{{-}}\NormalTok{ gamma\_i) }\OperatorTok{*}\NormalTok{ prob\_individual }\OperatorTok{+}\NormalTok{ gamma\_i }\OperatorTok{*}\NormalTok{ prob\_social}
\NormalTok{        prob\_final }\OperatorTok{=}\NormalTok{ jnp.where(bout\_i }\OperatorTok{\textgreater{}} \DecValTok{1}\NormalTok{, prob\_mix, prob\_individual)}

        \ControlFlowTok{return}\NormalTok{ AC.at[id\_i].}\BuiltInTok{set}\NormalTok{(ac\_new), prob\_final}

\NormalTok{    AC\_init }\OperatorTok{=}\NormalTok{ jnp.zeros((J, K))}
\NormalTok{    xs }\OperatorTok{=}\NormalTok{ (phi\_arr, gamma\_arr, fconf\_arr, Bpay\_arr, tech, }\BuiltInTok{id}\NormalTok{, bout, y\_prev, s, ps)}
\NormalTok{    \_, probs\_matrix }\OperatorTok{=}\NormalTok{ jax.lax.scan(step\_fn, AC\_init, xs)}
    
\NormalTok{    m.dist.categorical(}\StringTok{"obs"}\NormalTok{, probs}\OperatorTok{=}\NormalTok{probs\_matrix, obs}\OperatorTok{=}\NormalTok{tech)}

\NormalTok{m.fit(model)}
\NormalTok{m.summary()}
\end{Highlighting}
\end{Shaded}

\subsection{R}

\begin{Shaded}
\begin{Highlighting}[]
\FunctionTok{library}\NormalTok{(BayesianInference)}
\NormalTok{m }\OtherTok{\textless{}{-}} \FunctionTok{importBI}\NormalTok{(}\AttributeTok{platform=}\StringTok{\textquotesingle{}cpu\textquotesingle{}}\NormalTok{)}

\NormalTok{model }\OtherTok{\textless{}{-}} \ControlFlowTok{function}\NormalTok{(K, J, tech, id, bout, y\_prev, s, ps) \{}
  \CommentTok{\# Priors}
\NormalTok{  lam }\OtherTok{\textless{}{-}}\NormalTok{ m}\SpecialCharTok{$}\NormalTok{dist}\SpecialCharTok{$}\FunctionTok{exponential}\NormalTok{(}\AttributeTok{name=}\StringTok{"lambda"}\NormalTok{, }\AttributeTok{rate=}\FloatTok{1.0}\NormalTok{)}
\NormalTok{  mu }\OtherTok{\textless{}{-}}\NormalTok{ m}\SpecialCharTok{$}\NormalTok{dist}\SpecialCharTok{$}\FunctionTok{normal}\NormalTok{(}\AttributeTok{name=}\StringTok{"mu"}\NormalTok{, }\AttributeTok{loc=}\NormalTok{jnp}\SpecialCharTok{$}\FunctionTok{zeros}\NormalTok{(}\DecValTok{4}\NormalTok{L), }\AttributeTok{scale=}\FloatTok{1.0}\NormalTok{)}
\NormalTok{  sigma }\OtherTok{\textless{}{-}}\NormalTok{ m}\SpecialCharTok{$}\NormalTok{dist}\SpecialCharTok{$}\FunctionTok{exponential}\NormalTok{(}\AttributeTok{name=}\StringTok{"sigma"}\NormalTok{, }\AttributeTok{rate=}\NormalTok{jnp}\SpecialCharTok{$}\FunctionTok{full}\NormalTok{(}\DecValTok{4}\NormalTok{L, }\FloatTok{3.0}\NormalTok{))}
\NormalTok{  L\_Rho }\OtherTok{\textless{}{-}}\NormalTok{ m}\SpecialCharTok{$}\NormalTok{dist}\SpecialCharTok{$}\FunctionTok{lkj\_cholesky}\NormalTok{(}\AttributeTok{name=}\StringTok{"L\_Rho"}\NormalTok{, }\AttributeTok{dimension=}\DecValTok{4}\NormalTok{L, }\AttributeTok{concentration=}\FloatTok{3.0}\NormalTok{)}
\NormalTok{  z }\OtherTok{\textless{}{-}}\NormalTok{ m}\SpecialCharTok{$}\NormalTok{dist}\SpecialCharTok{$}\FunctionTok{normal}\NormalTok{(}\AttributeTok{name=}\StringTok{"z"}\NormalTok{, }\AttributeTok{loc=}\NormalTok{jnp}\SpecialCharTok{$}\FunctionTok{zeros}\NormalTok{(}\FunctionTok{tuple}\NormalTok{(}\DecValTok{4}\NormalTok{L, J)), }\AttributeTok{scale=}\FloatTok{1.0}\NormalTok{)}

  \CommentTok{\# ... Equivalent logic to Python implementation ...}
\NormalTok{\}}
\NormalTok{m}\SpecialCharTok{$}\FunctionTok{fit}\NormalTok{(model)}
\NormalTok{m}\SpecialCharTok{$}\FunctionTok{summary}\NormalTok{()}
\end{Highlighting}
\end{Shaded}

\subsection{Julia}

\begin{Shaded}
\begin{Highlighting}[]
\ImportTok{using} \BuiltInTok{BayesianInference}
\NormalTok{m }\OperatorTok{=} \FunctionTok{importBI}\NormalTok{(platform}\OperatorTok{=}\StringTok{"cpu"}\NormalTok{)}

\PreprocessorTok{@BI} \KeywordTok{function} \FunctionTok{model}\NormalTok{(K, J, tech, id, bout, y\_prev, s, ps)}
    \CommentTok{\# Priors}
\NormalTok{    lam }\OperatorTok{=}\NormalTok{ m.dist.}\FunctionTok{exponential}\NormalTok{(name}\OperatorTok{=}\StringTok{"lambda"}\NormalTok{, rate}\OperatorTok{=}\FloatTok{1.0}\NormalTok{)}
\NormalTok{    mu }\OperatorTok{=}\NormalTok{ m.dist.}\FunctionTok{normal}\NormalTok{(name}\OperatorTok{=}\StringTok{"mu"}\NormalTok{, loc}\OperatorTok{=}\NormalTok{jnp.}\FunctionTok{zeros}\NormalTok{(}\FloatTok{4}\NormalTok{), scale}\OperatorTok{=}\FloatTok{1.0}\NormalTok{)}
\NormalTok{    sigma }\OperatorTok{=}\NormalTok{ m.dist.}\FunctionTok{exponential}\NormalTok{(name}\OperatorTok{=}\StringTok{"sigma"}\NormalTok{, rate}\OperatorTok{=}\NormalTok{jnp.}\FunctionTok{full}\NormalTok{(}\FloatTok{4}\NormalTok{, }\FloatTok{3.0}\NormalTok{))}
\NormalTok{    L\_Rho }\OperatorTok{=}\NormalTok{ m.dist.}\FunctionTok{lkj\_cholesky}\NormalTok{(name}\OperatorTok{=}\StringTok{"L\_Rho"}\NormalTok{, dimension}\OperatorTok{=}\FloatTok{4}\NormalTok{, concentration}\OperatorTok{=}\FloatTok{3.0}\NormalTok{)}
\NormalTok{    z }\OperatorTok{=}\NormalTok{ m.dist.}\FunctionTok{normal}\NormalTok{(name}\OperatorTok{=}\StringTok{"z"}\NormalTok{, loc}\OperatorTok{=}\NormalTok{jnp.}\FunctionTok{zeros}\NormalTok{(}\FloatTok{4}\NormalTok{, J), scale}\OperatorTok{=}\FloatTok{1.0}\NormalTok{)}

    \CommentTok{\# ... Equivalent logic to Python implementation ...}
\KeywordTok{end}
\NormalTok{m.}\FunctionTok{fit}\NormalTok{(model)}
\NormalTok{m.}\FunctionTok{summary}\NormalTok{()}
\end{Highlighting}
\end{Shaded}

\subsubsection{Mathematical Details}\label{mathematical-details}

The ``Experience-Weighted Attraction'' (EWA) model links individual
learning rates and social frequency to behavioral choice. Throughout
this section, we use a consistent three-part index: \(i\) denotes the
individual, \(k \in \{1, \ldots, K\}\) denotes the technique, and \(t\)
denotes the observation (time step).

\textbf{1. Mixture Likelihood}

At each observation, the probability of choosing a single technique from
\(K\) options is expressed through a \emph{Categorical} likelihood: \[
Y_{[i,t]} \sim \text{Categorical}(\boldsymbol{\theta}_{[i,t]})
\] where:

\begin{itemize}
\tightlist
\item
  \(Y_{[i,t]}\) is the technique chosen by individual \(i\) at time step
  \(t\).
\item
  \(\boldsymbol{\theta}_{[i,t]}\) is a probability vector of length
  \(K\) giving the probability of each technique for individual \(i\) at
  time step \(t\).
\end{itemize}

Unlike a standard Categorical model where \(\boldsymbol{\theta}\) is
fixed for a group of observations, here \(\boldsymbol{\theta}_{[i,t]}\)
is a \textbf{different vector at every time step} for every individual,
because the attraction scores and social cues update sequentially.

The probability that individual \(i\) chooses technique \(k\) at time
\(t\) decomposes as: \[
\theta_{[i,k,t]} = (1 - \gamma_i)\, I_{[i,k,t]} + \gamma_i\, S_{[i,k,t]}
\] \[
\gamma_i = \text{sigmoid}(\mu_2 + a_{id}[i,2])
\]

Where:

\begin{itemize}
\tightlist
\item
  \(I_{[i,k,t]}\) is the asocial choice probability: derived from
  individual \(i\)'s personal experience with technique \(k\).
\item
  \(S_{[i,k,t]}\) is the social choice probability: derived from
  observing conspecifics use technique \(k\).
\item
  \(\gamma_i \in [0, 1]\) is individual \(i\)'s social learning weight:
  at \(0\) the individual relies entirely on personal experience; at
  \(1\) entirely on social information. It is derived from the
  hierarchical prior (see §2):
\end{itemize}

\textbf{2. Asocial Component} --- \(I_{[i,k,t]}\)

The asocial probability is a softmax over accumulated attraction scores
--- it converts raw attraction values into a proper probability
distribution over techniques: \[
I_{[i,k,t]} = \frac{\exp(\lambda \cdot A_{[i,k,t]})}{\sum_{k'=1}^{K} \exp(\lambda \cdot A_{[i,k',t]})}
\] \[
\lambda \sim \text{Exponential}(1)
\]

The numerator scores technique \(k\) by the individual's scaled
attraction to it. The denominator sums this over all \(K\) techniques
(\(k' = 1,\ldots,K\)), ensuring the probabilities sum to 1 and that each
technique is evaluated relative to all alternatives.

Where:

\begin{itemize}
\tightlist
\item
  \(A_{[i,k,t]}\) is the attraction score of individual \(i\) for
  technique \(k\) at time \(t\). It is a running exponentially weighted
  average of the personal payoffs individual \(i\) has obtained from
  technique \(k\) across past bouts. Larger values indicate a stronger
  learned payoff expectation.
\item
  \(A_{[i,k',t]}\) in the denominator ranges over the attraction scores
  for all \(K\) techniques, making choice proportional to each option's
  scaled attraction relative to the full set.
\item
  \(\lambda > 0\) is the multinomial sensitivity (inverse temperature):
  higher values make choice more deterministic toward the
  highest-attraction option; \(\lambda \to 0\) yields random choice. The
  Exponential(1) prior keeps \(\lambda\) positive and places most mass
  near zero, expressing a prior expectation that choice is moderately
  stochastic. Larger values are possible but require evidence from the
  data.
\end{itemize}

\emph{Attraction scores update between foraging bouts via the EWA rule:}
\[
A_{[i,k,t+1]} = (1 - \phi_i)\, A_{[i,k,t]} + \phi_i\, p_{[i,k,t]}
\]

\[
\phi_i = \text{sigmoid}(\mu_1 + a_{id}[i,1])
\]

\[
\boldsymbol{\mu} \sim \mathcal{N}(\mathbf{0}_4,\, 1)
\] \[
\boldsymbol{\sigma} \sim \text{Exponential}(3)
\] \[
\mathbf{L}_\rho \sim \text{LKJCholesky}(4,\, \eta = 3)
\]

\[
\mathbf{Z} \sim \mathcal{N}(\mathbf{0}_{4 \times J},\, 1)
\] \[
\mathbf{a}_{id} = \bigl(\mathbf{L}_\rho \cdot \operatorname{diag}(\boldsymbol{\sigma}) \cdot \mathbf{Z}\bigr)^\top
\]

Where:

\begin{itemize}
\tightlist
\item
  \(A_{[i,k,t]}\) is the attraction score of individual \(i\) for
  technique \(k\) at time \(t\).
\item
  \(p_{[i,k,t]}\) is the personal yield (payoff) individual \(i\)
  obtained from technique \(k\) at time \(t\) --- concretely, a measure
  of food-extraction efficiency (e.g.~rate of fruit opening).
\item
  \(\phi_i \in [0, 1]\) is individual \(i\)'s attraction updating
  weight. Near \(1\): attractions closely track the most recent payoff.
  Near \(0\): strong memory of cumulative past experience. It is
  constructed from a multivariate hierarchical prior shared by all four
  individual-level learning parameters (\(\phi_i\), \(\gamma_i\),
  \(f_i\), \(\beta_{i}\))
\item
  \(\boldsymbol{\mu} \in \mathbb{R}^4\) are population-level means
\item
  \(a_{id}[i, \cdot]\) are individual-level correlated deviations
  reconstructed via non-centered parameterization
\item
  \(\boldsymbol{\sigma} \in \mathbb{R}^4_+\) are the per-parameter
  standard deviations (shrinkage parameters)
\item
  \(\mathbf{L}_\rho\) is the Cholesky factor of the \(4\times4\)
  correlation matrix: the LKJ prior (\(\eta=3\)) moderately shrinks
  correlations toward zero.
\item
  \(\mathbf{Z} \in \mathbb{R}^{J \times 4}\) is a matrix of standard
  normal deviates (one row per individual).
\end{itemize}

\textbf{3. Social Component} --- \(S_{[i,k,t]}\)

The social probability combines payoff bias with frequency-dependent
conformity: \[
S_{[i,k,t]} = \frac{N_{[i,k,t]}^{f_i}\, \exp(\beta_{i}\, \pi_{[i,k,t]})}{\sum_{k'=1}^{K} N_{[i,k',t]}^{f_i}\, \exp(\beta_{i}\, \pi_{[i,k',t]})}
\] \[
f_i = \exp(\mu_3 + a_{id}[i,3])
\] \[
\beta_{i} = \mu_4 + a_{id}[i,4]
\]

Where:

\begin{itemize}
\tightlist
\item
  \(N_{[i,k,t]}\) is the \emph{frequency-bias cue}: the number of times
  individual \(i\) observed conspecifics performing technique \(k\) at
  time \(t\).
\item
  \(f_i > 0\) is individual \(i\)'s conformity exponent. \(f_i > 1\):
  positive frequency-dependent bias (disproportionate copying of the
  most common technique); \(f_i = 1\): purely linear frequency
  weighting; \(f_i < 1\): tendency to avoid the majority. The
  exponential link ensures \(f_i > 0\); the prior
  \(\mu_3 \sim \mathcal{N}(0,1)\) corresponds to \(f \approx 1\) a
  priori.
\item
  \(N_{[i,k',t]}^{f_i}\) in the denominator ranges over all \(K\)
  techniques with the same conformity exponent \(f_i\), normalising the
  social probabilities.
\item
  \(\pi_{[i,k,t]}\) is the payoff observed from social demonstrators
  using technique \(k\) at time \(t\).
\item
  \(\beta_{i}\) is individual \(i\)'s payoff-bias coefficient: positive
  values mean the individual preferentially copies techniques
  demonstrated with higher payoff. It is on the unconstrained real
  scale; the prior \(\mu_4 \sim \mathcal{N}(0,1)\) expresses prior
  uncertainty about the direction of payoff bias.
\end{itemize}

The numerator scores technique \(k\) by two multiplicative components:
how often it was observed being used by conspecifics
(\(N_{[i,k,t]}^{f_i}\), frequency bias) and how profitable it appeared
in their hands (\(\exp(\beta_{i}\, \pi_{[i,k,t]})\), payoff bias). The
denominator (\(k' = 1, \ldots, K\)) sums these scores over all
techniques to normalise to a probability. Together they implement the
idea that individuals copy techniques that are simultaneously common
\emph{and} profitable.

\begin{center}\rule{0.5\linewidth}{0.5pt}\end{center}

\subsection{2. Social Learning (EWA-ILV)}\label{social-learning-ewa-ilv}

The ``EWA-ILV'' model expands the EWA version by incorporating a wider
array of social biases (kinship, prestige, age similarity, etc.) and
modeling the linear influence of individual age on the learning
parameters \(\phi\) and \(\gamma\).

\subsubsection{Example}\label{example-1}

\subsection{Python}

\begin{Shaded}
\begin{Highlighting}[]
\ImportTok{from}\NormalTok{ BI }\ImportTok{import}\NormalTok{ bi}
\ImportTok{import}\NormalTok{ jax.numpy }\ImportTok{as}\NormalTok{ jnp}
\ImportTok{import}\NormalTok{ jax}

\NormalTok{m }\OperatorTok{=}\NormalTok{ bi(platform}\OperatorTok{=}\StringTok{\textquotesingle{}cpu\textquotesingle{}}\NormalTok{)}

\KeywordTok{def}\NormalTok{ model(K, J, tech, }\BuiltInTok{id}\NormalTok{, bout, y\_prev, s, ps, ks, pr, co, yo, age):}
    \CommentTok{\# Priors (8 varying effects)}
\NormalTok{    lam   }\OperatorTok{=}\NormalTok{ m.dist.exponential(name}\OperatorTok{=}\StringTok{"lambda"}\NormalTok{, rate}\OperatorTok{=}\FloatTok{1.0}\NormalTok{)}
\NormalTok{    mu    }\OperatorTok{=}\NormalTok{ m.dist.normal(name}\OperatorTok{=}\StringTok{"mu"}\NormalTok{, loc}\OperatorTok{=}\NormalTok{jnp.zeros(}\DecValTok{8}\NormalTok{), scale}\OperatorTok{=}\FloatTok{1.0}\NormalTok{)}
\NormalTok{    sigma }\OperatorTok{=}\NormalTok{ m.dist.exponential(name}\OperatorTok{=}\StringTok{"sigma"}\NormalTok{, rate}\OperatorTok{=}\NormalTok{jnp.full(}\DecValTok{8}\NormalTok{, }\FloatTok{3.0}\NormalTok{))}
\NormalTok{    L\_Rho }\OperatorTok{=}\NormalTok{ m.dist.lkj\_cholesky(name}\OperatorTok{=}\StringTok{"L\_Rho"}\NormalTok{, dimension}\OperatorTok{=}\DecValTok{8}\NormalTok{, concentration}\OperatorTok{=}\FloatTok{3.0}\NormalTok{)}
\NormalTok{    z     }\OperatorTok{=}\NormalTok{ m.dist.normal(name}\OperatorTok{=}\StringTok{"z"}\NormalTok{, loc}\OperatorTok{=}\NormalTok{jnp.zeros((}\DecValTok{8}\NormalTok{, J)), scale}\OperatorTok{=}\FloatTok{1.0}\NormalTok{)}
\NormalTok{    b\_age }\OperatorTok{=}\NormalTok{ m.dist.normal(name}\OperatorTok{=}\StringTok{"b\_age"}\NormalTok{, loc}\OperatorTok{=}\NormalTok{jnp.zeros(}\DecValTok{2}\NormalTok{), scale}\OperatorTok{=}\FloatTok{1.0}\NormalTok{)}

    \CommentTok{\# Varying effects}
\NormalTok{    L\_scaled }\OperatorTok{=}\NormalTok{ L\_Rho }\OperatorTok{*}\NormalTok{ sigma[:, }\VariableTok{None}\NormalTok{]}
\NormalTok{    a\_id     }\OperatorTok{=}\NormalTok{ (L\_scaled }\OperatorTok{@}\NormalTok{ z).T}

    \CommentTok{\# Parameters per individual with age effects}
\NormalTok{    phi\_arr   }\OperatorTok{=}\NormalTok{ jax.nn.sigmoid(mu[}\DecValTok{0}\NormalTok{] }\OperatorTok{+}\NormalTok{ a\_id[}\BuiltInTok{id}\NormalTok{, }\DecValTok{0}\NormalTok{] }\OperatorTok{+}\NormalTok{ b\_age[}\DecValTok{0}\NormalTok{] }\OperatorTok{*}\NormalTok{ age)}
\NormalTok{    gamma\_arr }\OperatorTok{=}\NormalTok{ jax.nn.sigmoid(mu[}\DecValTok{1}\NormalTok{] }\OperatorTok{+}\NormalTok{ a\_id[}\BuiltInTok{id}\NormalTok{, }\DecValTok{1}\NormalTok{] }\OperatorTok{+}\NormalTok{ b\_age[}\DecValTok{1}\NormalTok{] }\OperatorTok{*}\NormalTok{ age)}
\NormalTok{    fconf\_arr }\OperatorTok{=}\NormalTok{ jnp.exp(mu[}\DecValTok{2}\NormalTok{] }\OperatorTok{+}\NormalTok{ a\_id[}\BuiltInTok{id}\NormalTok{, }\DecValTok{2}\NormalTok{])}
\NormalTok{    B\_pay     }\OperatorTok{=}\NormalTok{ mu[}\DecValTok{3}\NormalTok{] }\OperatorTok{+}\NormalTok{ a\_id[}\BuiltInTok{id}\NormalTok{, }\DecValTok{3}\NormalTok{]}
\NormalTok{    B\_kin     }\OperatorTok{=}\NormalTok{ mu[}\DecValTok{4}\NormalTok{] }\OperatorTok{+}\NormalTok{ a\_id[}\BuiltInTok{id}\NormalTok{, }\DecValTok{4}\NormalTok{]}
\NormalTok{    B\_pres    }\OperatorTok{=}\NormalTok{ mu[}\DecValTok{5}\NormalTok{] }\OperatorTok{+}\NormalTok{ a\_id[}\BuiltInTok{id}\NormalTok{, }\DecValTok{5}\NormalTok{]}
\NormalTok{    B\_coho    }\OperatorTok{=}\NormalTok{ mu[}\DecValTok{6}\NormalTok{] }\OperatorTok{+}\NormalTok{ a\_id[}\BuiltInTok{id}\NormalTok{, }\DecValTok{6}\NormalTok{]}
\NormalTok{    B\_yob     }\OperatorTok{=}\NormalTok{ mu[}\DecValTok{7}\NormalTok{] }\OperatorTok{+}\NormalTok{ a\_id[}\BuiltInTok{id}\NormalTok{, }\DecValTok{7}\NormalTok{]}

    \CommentTok{\# Scan carry: Attractions}
    \KeywordTok{def}\NormalTok{ step\_fn(carry, x):}
\NormalTok{        AC }\OperatorTok{=}\NormalTok{ carry}
\NormalTok{        (phi\_i, gamma\_i, fconf\_i, Bp, Bk, Bpr, Bc, By, t\_i, id\_i, b\_i, py\_i, s\_i, ps\_i, ks\_i, pr\_i, co\_i, yo\_i) }\OperatorTok{=}\NormalTok{ x}

\NormalTok{        ac\_new }\OperatorTok{=}\NormalTok{ jnp.where(b\_i }\OperatorTok{==} \DecValTok{1}\NormalTok{, jnp.zeros(K), (}\FloatTok{1.0} \OperatorTok{{-}}\NormalTok{ phi\_i) }\OperatorTok{*}\NormalTok{ AC[id\_i] }\OperatorTok{+}\NormalTok{ phi\_i }\OperatorTok{*}\NormalTok{ py\_i)}
        
        \CommentTok{\# Asocial choice probability}
\NormalTok{        prob\_individual }\OperatorTok{=}\NormalTok{ jax.nn.softmax(lam }\OperatorTok{*}\NormalTok{ ac\_new)}

        \CommentTok{\# Multi{-}cue social influence}
\NormalTok{        lin\_rest }\OperatorTok{=}\NormalTok{ jnp.exp(Bp }\OperatorTok{*}\NormalTok{ ps\_i[}\DecValTok{1}\NormalTok{:] }\OperatorTok{+}\NormalTok{ Bk }\OperatorTok{*}\NormalTok{ ks\_i[}\DecValTok{1}\NormalTok{:] }\OperatorTok{+}\NormalTok{ Bpr }\OperatorTok{*}\NormalTok{ pr\_i[}\DecValTok{1}\NormalTok{:] }\OperatorTok{+}\NormalTok{ Bc }\OperatorTok{*}\NormalTok{ co\_i[}\DecValTok{1}\NormalTok{:] }\OperatorTok{+}\NormalTok{ By }\OperatorTok{*}\NormalTok{ yo\_i[}\DecValTok{1}\NormalTok{:])}
\NormalTok{        lin\_mod }\OperatorTok{=}\NormalTok{ jnp.concatenate([jnp.ones(}\DecValTok{1}\NormalTok{), lin\_rest]) }\OperatorTok{*}\NormalTok{ jnp.power(s\_i, fconf\_i)}
\NormalTok{        safe\_denom }\OperatorTok{=}\NormalTok{ jnp.where(jnp.}\BuiltInTok{sum}\NormalTok{(lin\_mod) }\OperatorTok{\textgreater{}} \DecValTok{0}\NormalTok{, jnp.}\BuiltInTok{sum}\NormalTok{(lin\_mod), }\FloatTok{1.0}\NormalTok{)}
\NormalTok{        prob\_social }\OperatorTok{=}\NormalTok{ lin\_mod }\OperatorTok{/}\NormalTok{ safe\_denom}

        \CommentTok{\# Mixture probability}
\NormalTok{        prob\_mix }\OperatorTok{=}\NormalTok{ (}\FloatTok{1.0} \OperatorTok{{-}}\NormalTok{ gamma\_i) }\OperatorTok{*}\NormalTok{ prob\_individual }\OperatorTok{+}\NormalTok{ gamma\_i }\OperatorTok{*}\NormalTok{ prob\_social}
\NormalTok{        has\_social\_info }\OperatorTok{=}\NormalTok{ jnp.logical\_and(b\_i }\OperatorTok{\textgreater{}} \DecValTok{1}\NormalTok{, jnp.}\BuiltInTok{sum}\NormalTok{(s\_i) }\OperatorTok{\textgreater{}} \DecValTok{0}\NormalTok{)}
\NormalTok{        prob\_final }\OperatorTok{=}\NormalTok{ jnp.where(has\_social\_info, prob\_mix, prob\_individual)}

        \ControlFlowTok{return}\NormalTok{ AC.at[id\_i].}\BuiltInTok{set}\NormalTok{(ac\_new), prob\_final}

\NormalTok{    AC\_init }\OperatorTok{=}\NormalTok{ jnp.zeros((J, K))}
\NormalTok{    xs }\OperatorTok{=}\NormalTok{ (phi\_arr, gamma\_arr, fconf\_arr, B\_pay, B\_kin, B\_pres, B\_coho, B\_yob,}
\NormalTok{          tech, }\BuiltInTok{id}\NormalTok{, bout, y\_prev, s, ps, ks, pr, co, yo)}
\NormalTok{    \_, probs\_matrix }\OperatorTok{=}\NormalTok{ jax.lax.scan(step\_fn, AC\_init, xs)}
    
\NormalTok{    m.dist.categorical(}\StringTok{"obs"}\NormalTok{, probs}\OperatorTok{=}\NormalTok{probs\_matrix, obs}\OperatorTok{=}\NormalTok{tech)}

\NormalTok{m.fit(model)}
\NormalTok{m.summary()}
\end{Highlighting}
\end{Shaded}

\subsection{R}

\begin{Shaded}
\begin{Highlighting}[]
\CommentTok{\# Equivalent R implementation with 8 effects and age slopes...}
\end{Highlighting}
\end{Shaded}

\subsection{Julia}

\begin{Shaded}
\begin{Highlighting}[]
\CommentTok{\# Equivalent Julia implementation with 8 effects and age slopes...}
\end{Highlighting}
\end{Shaded}

\subsubsection{Mathematical Details}\label{mathematical-details-1}

The ``Experience-Weighted Attraction with Individual Level Covariates''
(EWA-ILV) model extends the EWA model in two ways: it incorporates a
wider array of social cues in the social learning component, and it
allows an individual's age to shift key learning parameters.

\textbf{1. Mixture Likelihood}

The observation likelihood has the same mixture form as the EWA model,
using a three-part index: individual \(i\), technique \(k\), and time
step \(t\): \[
\theta_{[i,k,t]} = (1 - \gamma_i)\, I_{[i,k,t]} + \gamma_i\, S_{[i,k,t]}
\]

\textbf{2. Social Component} --- \(S_{[i,k,t]}\)

The social probability combines frequency-dependent conformity with a
multi-cue log-linear model: \[
S_{[i,k,t]} = \frac{N_{[i,k,t]}^{f_i}\, \exp(B_{[i,k,t]})}{\sum_{k'=1}^{K} N_{[i,k',t]}^{f_i}\, \exp(B_{[i,k',t]})}
\] Where \(B_{[i,1,t]} = 0\) by convention (technique 1 is the aliased
baseline), and for \(k > 1\): \[
B_{[i,k,t]} = \beta_{\text{pay},i}\, \pi_{[i,k,t]} + \beta_{\text{kin},i}\, \kappa_{[i,k,t]} + \beta_{\text{rank},i}\, \zeta_{[i,k,t]} + \beta_{\text{cohort},i}\, \eta_{[i,k,t]} + \beta_{\text{age},i}\, \alpha_{[i,k,t]}
\]

\textbf{3. Individual-Level Parameters, Age Effects, and Priors}

All eight learning parameters are drawn from a correlated multivariate
hierarchy. For each individual \(i\), the varying parameters and
deviations are constructed as: \[
\phi_i = \text{sigmoid}\bigl(\mu_0 + a_{id}[i, 0] + \beta_{\phi} \cdot \text{age}_i\bigr)
\] \[
\gamma_i = \text{sigmoid}\bigl(\mu_1 + a_{id}[i, 1] + \beta_{\gamma} \cdot \text{age}_i\bigr)
\] \[
f_i = \exp(\mu_2 + a_{id}[i, 2])
\] \[
\beta_{\text{pay},i} = \mu_3 + a_{id}[i, 3]
\] \[
\beta_{\text{kin},i} = \mu_4 + a_{id}[i, 4]
\] \[
\beta_{\text{rank},i} = \mu_5 + a_{id}[i, 5]
\] \[
\beta_{\text{cohort},i} = \mu_6 + a_{id}[i, 6]
\] \[
\beta_{\text{age},i} = \mu_7 + a_{id}[i, 7]
\] \[
\mathbf{a}_{id} = (\mathbf{L}_{\rho} \operatorname{diag}(\boldsymbol{\sigma}) \, \mathbf{Z})^\top
\] \[
\boldsymbol{\mu} \sim \mathcal{N}(\mathbf{0}_8, 1)
\] \[
\boldsymbol{\sigma} \sim \text{Exponential}(\text{rate}=3)
\] \[
\mathbf{L}_{\rho} \sim \text{LKJCholesky}(8, \eta=3)
\] \[
\mathbf{Z} \sim \mathcal{N}(\mathbf{0}_{8 \times J}, 1)
\] \[
\beta_{\phi} \sim \mathcal{N}(0, 1), \quad \beta_{\gamma} \sim \mathcal{N}(0, 1)
\] \[
\lambda \sim \text{Exponential}(1)
\]

Where:

\begin{itemize}
\tightlist
\item
  \(N_{[i,k,t]}\) is the frequency cue: the number of times individual
  \(i\) observed conspecifics performing technique \(k\) at time \(t\).
\item
  \(f_i > 0\) is individual \(i\)'s conformity exponent, constructed
  with the exponential link to guarantee positive values. The population
  prior mean \(\mu_2 \sim \mathcal{N}(0,1)\) corresponds to
  \(f \approx 1\) (linear frequency copying) a priori.
\item
  \(\pi_{[i,k,t]}\) is the payoff observed from social demonstrators
  using technique \(k\) at time \(t\). The prior
  \(\mu_3 \sim \mathcal{N}(0,1)\) expresses uncertainty on payoff bias
  direction and magnitude.
\item
  \(\kappa_{[i,k,t]}\) is the kin bias cue: whether the demonstrators
  using technique \(k\) at time \(t\) are matrilineal kin of individual
  \(i\). The coefficient \(\beta_{\text{kin},i}\) models the kinship
  bias, with population prior mean \(\mu_4 \sim \mathcal{N}(0,1)\).
\item
  \(\zeta_{[i,k,t]}\) is the rank/prestige cue: whether the
  demonstrators using technique \(k\) at time \(t\) are high-ranking
  (e.g., alpha) individuals. The coefficient \(\beta_{\text{rank},i}\)
  models rank bias, with population prior mean
  \(\mu_5 \sim \mathcal{N}(0,1)\).
\item
  \(\eta_{[i,k,t]}\) is the cohort cue: the age-similarity between
  individual \(i\) and demonstrators using technique \(k\) at time
  \(t\). The coefficient \(\beta_{\text{cohort},i}\) models cohort bias,
  with population prior mean \(\mu_6 \sim \mathcal{N}(0,1)\).
\item
  \(\alpha_{[i,k,t]}\) is the age-bias cue: a year-of-birth prestige
  bias towards older/more experienced demonstrators using technique
  \(k\) at time \(t\). The coefficient \(\beta_{\text{age},i}\) models
  age-based prestige bias, with population prior mean
  \(\mu_7 \sim \mathcal{N}(0,1)\).
\item
  \(\phi_i \in [0, 1]\) and \(\gamma_i \in [0, 1]\) are individual-level
  attraction-updating and social-learning weights, constructed with the
  logistic/sigmoid link.
\item
  \(\text{age}_i\) is the centered age of individual \(i\).
\item
  \(\beta_{\phi}\) and \(\beta_{\gamma}\) are the population-level age
  slopes on the log-odds scale, determining how an individual's age
  alters their attraction updating and social learning weight.
\item
  \(\mathbf{a}_{id} \in \mathbb{R}^{J \times 8}\) is the matrix of
  individual-level deviations across the 8 varying parameters,
  reconstructed using the non-centered parameterization to optimize
  Hamiltonian Monte Carlo performance.
\item
  \(\mathbf{L}_{\rho} \in \mathbb{R}^{8 \times 8}\) is the Cholesky
  factor of the correlation matrix, mapping correlations among all 8
  parameters.
\item
  \(\boldsymbol{\sigma} \in \mathbb{R}^8_+\) represents the standard
  deviations (shrinkage parameters) for each of the 8 parameters.
\item
  \(\mathbf{Z} \in \mathbb{R}^{8 \times J}\) is a matrix of standard
  normal deviates.
\item
  \(\lambda > 0\) is the multinomial sensitivity parameter, governing
  choice deterministic behavior.
\end{itemize}

\begin{center}\rule{0.5\linewidth}{0.5pt}\end{center}

\subsection{Notes}\label{notes}

\begin{tcolorbox}[enhanced jigsaw, colback=white, toprule=.15mm, coltitle=black, colframe=quarto-callout-note-color-frame, breakable, opacityback=0, bottomrule=.15mm, leftrule=.75mm, opacitybacktitle=0.6, colbacktitle=quarto-callout-note-color!10!white, title=\textcolor{quarto-callout-note-color}{\faInfo}\hspace{0.5em}{Note}, left=2mm, bottomtitle=1mm, toptitle=1mm, titlerule=0mm, arc=.35mm, rightrule=.15mm]

\begin{itemize}
\tightlist
\item
  The \texttt{EWA-ILV} model is significantly more complex and requires
  larger sample sizes to resolve the correlations between varying
  effects.
\item
  Cholesky decomposition of the correlation matrix is used for both
  models to ensure positive-definiteness.
\end{itemize}

\end{tcolorbox}

\subsection*{Reference(s)}\label{references}
\addcontentsline{toc}{subsection}{Reference(s)}

\phantomsection\label{refs}
\begin{CSLReferences}{1}{0}
\bibitem[\citeproctext]{ref-Barrett2017}
Barrett, Brendan J., Richard L. McElreath, and Susan E. Perry. 2017.
{``Pay-Off-Biased Social Learning Underlies the Diffusion of Novel
Extractive Foraging Traditions in a Wild Primate.''} \emph{Proceedings
of the Royal Society B: Biological Sciences} 284 (1856): 20170358.
\url{https://doi.org/10.1098/rspb.2017.0358}.

\end{CSLReferences}




\end{document}
